\documentclass{article}
\usepackage{../assignment_preamble}

\title{Weekly Problem 4}
\author{Ravi Kini}
\date{April 27, 2026}

\begin{document}

\maketitle

\problem
\subproblem{(a)}
\begin{equation*}
\begin{nd}[rules=paradox,justformat=just]
\hypo {1} {\lnot p \lor q}
\open
\hypo {2} {p}
\open
\hypo {3} {\lnot p}
\have {4} {\lnot p \lor (p \land q)} \oi{3}
\close
\open
\hypo {5} {q}
\have {6} {p \land q} \ai{2,5}
\have {7} {\lnot p \lor (p \land q)} \oi{6}
\close
\have {8} {\lnot p \lor (p \land q)} \oe{1,2-4,5-7}
\close
\open
\hypo {9} {\lnot p}
\have {10} {\lnot p \lor (p \land q)} \oe{9}
\close
\have {11} {\lnot p \lor (p \land q)} \tnd{2-8,9-10}
\have {12} {p \Rightarrow (p \land q)} \mci{11}
\end{nd}
\end{equation*}
We then see that $\lnot p \lor q \vdash p \Rightarrow (p \land q)$ is a valid argument in the Logic of Paradox.

\subproblem{(b)}
Let $p = i$, $q = 1$, and $r = 0$. We then see that $p \land q = i$, so $(p \land q) \Rightarrow r = i$. However, $q \Rightarrow r = 0$. Consequently, the argument $p, (p \land q) \Rightarrow r \vdash q \Rightarrow r$ is an invalid argument in the Logic of Paradox.

\end{document}